\documentclass{article}
\usepackage{pgfplots}
\pgfplotsset{compat=1.18}
\usepackage{filecontents}

% Points d'entraînement du réseau
\begin{filecontents}{points.dat}
x y z
3 5 4.712
5 1 0.906
10 2 1.681
\end{filecontents}

% Quelques points supplémentaires calculés par le réseau
\begin{filecontents}{predictions.dat}
x y z
1 1 0.8
2 2 1.5
4 3 2.2
6 4 3.0
8 5 3.8
9 1 0.4
7 2 1.1
6 3 1.8
5 4 2.5
4 5 3.2
2 4 2.8
3 1 0.7
8 3 1.9
7 4 3.1
9 5 4.0
\end{filecontents}

\begin{document}

\begin{figure}[htbp]
\centering

% Version 3D simple - juste des points
\begin{tikzpicture}
\begin{axis}[
    width=12cm,
    height=9cm,
    xlabel={$x_1$},
    ylabel={$x_2$},
    zlabel={$\hat{y}$},
    title={Prédictions du réseau de neurones (architecture 2→3→1)},
    view={45}{35},
    grid=major,
    legend pos=north east,
]

% Points d'entraînement originaux (gros points rouges)
\addplot3[
    only marks,
    mark=*,
    mark size=5pt,
    color=red,
] table {points.dat};

% Prédictions supplémentaires (petits points bleus)  
\addplot3[
    only marks,
    mark=o,
    mark size=3pt,
    color=blue,
] table {predictions.dat};

\legend{Points d'entraînement, Autres prédictions}

\end{axis}
\end{tikzpicture}

\caption{Visualisation 3D des prédictions du réseau de neurones. Les points rouges sont les données d'entraînement originales $(3,5) \rightarrow 4.71$, $(5,1) \rightarrow 0.91$, $(10,2) \rightarrow 1.68$. Les points bleus montrent d'autres prédictions du réseau pour illustrer la fonction apprise.}

\end{figure}

\end{document}